% Layout of the print edition, loaded by tc.tex when \tcprint is
% defined (see tc-print.tex).
%
% Trim size 13.5 x 21.5 cm, a standard BoD format. The text block is
% 100 mm wide: wide enough for 10 pt body text at about 65 characters
% a line, and for 67 columns of code at 8 pt, which is the limit that
% `make check` enforces on code blocks.

\usepackage[paperwidth=135mm,paperheight=215mm,
            inner=20mm,outer=15mm,top=18mm,bottom=20mm,
            headheight=12pt,headsep=5mm,footskip=9mm]{geometry}

% Code and diagrams at 8 pt.
\RecustomVerbatimEnvironment{Highlighting}{Verbatim}%
  {commandchars=\\\{\},fontsize=\footnotesize}
\RecustomVerbatimEnvironment{verbatim}{Verbatim}{fontsize=\footnotesize}

% Running heads: page number outside, chapter on the left page, section
% on the right. Chapter openings show only a page number at the foot.
\usepackage{fancyhdr}
\pagestyle{fancy}
\fancyhf{}
\fancyhead[LE,RO]{\small\thepage}
\fancyhead[RE]{\small\nouppercase{\leftmark}}
\fancyhead[LO]{\small\nouppercase{\rightmark}}
\renewcommand{\headrulewidth}{0pt}
\fancypagestyle{plain}{%
  \fancyhf{}%
  \fancyfoot[C]{\small\thepage}%
  \renewcommand{\headrulewidth}{0pt}}

% Blank pages left by chapters opening on the right carry no heading.
\makeatletter
\def\cleardoublepage{\clearpage\if@twoside\ifodd\c@page\else
  \hbox{}\thispagestyle{empty}\newpage\fi\fi}
\makeatother

% Tables at 9 pt with closer columns, so that the widest fit the
% 100 mm text block.
\AtBeginEnvironment{longtable}{\small}
\setlength{\tabcolsep}{3pt}

% Black only. The screen edition colours its code; in print, colour
% would turn those pages into colour pages, which are printed and
% priced separately. The highlighting macros set their colour with
% \textcolor[model]{spec}{text}, so here it keeps the text and drops
% the colour, with or without the optional colour model.
\renewcommand{\textcolor}[3][]{#3}

% A thread-sewn binding is made of folded sheets of four pages, so the
% total page count must be a multiple of four. The printer's order also
% fixes the exact count, \tcpages in tc-print.tex. Blank pages are added
% at the end to reach it; a book that has grown past it is padded only
% to a multiple of four, and the log says so (make print warns too).
\makeatletter
\newcommand{\tcpadto}[1]{%
  \ifnum\ReadonlyShipoutCounter<#1\relax
    \null\thispagestyle{empty}\newpage
    \expandafter\tcpadto\expandafter{\number#1}%
  \fi}
\newcommand{\tcpadfour}{%
  \count@=\ReadonlyShipoutCounter
  \divide\count@ 4 \multiply\count@ 4
  \ifnum\count@=\ReadonlyShipoutCounter\else
    \null\thispagestyle{empty}\newpage
    \expandafter\tcpadfour
  \fi}
\newcommand{\tcpad}{%
  \ifdefined\tcpages
    \ifnum\ReadonlyShipoutCounter>\tcpages\relax
      \typeout{tc-print: \the\ReadonlyShipoutCounter\space pages, more
        than the \tcpages\space ordered}%
    \else
      \tcpadto{\tcpages}%
    \fi
  \fi
  \tcpadfour}
\makeatother
\AtEndDocument{\clearpage\tcpad}

% Contents: room for two-digit section numbers such as 12.10, which
% the book class's default widths run into the title.
\makeatletter
\renewcommand*\l@section{\@dottedtocline{1}{1.5em}{3.0em}}
\renewcommand*\l@subsection{\@dottedtocline{2}{4.5em}{3.8em}}
\makeatother

% Keep a heading with at least a few lines of what follows it, so that
% a heading is never left alone at the foot of a page.
\usepackage{needspace}
\pretocmd{\section}{\needspace{8\baselineskip}}{}{}
\pretocmd{\subsection}{\needspace{7\baselineskip}}{}{}
\pretocmd{\subsubsection}{\needspace{6\baselineskip}}{}{}
